\scanpage{69}\mspage{60}

des $|Y|$, $Y\in\widetilde X$, est une figure élémentaire (de support $|X|$, bien sûr), car par (i) les $|Y|\subset |X|$.)\unsure{,|.} Cela résulte en effet du critère (\unsure{ivi|iVi|iii}) de la prop. p. 37. (\uline{NB} on n'a utilisé (i) que \unsure{pour} le cas $X\leq Y$.) Il est clair que cette application $f$ est croissante pour $\leq$,
% slanted
\margin{condition d) p. 28}
\deleted{et que} \inserted{et qu'elle \unsure{induit}, $\forall X\in M$, $\widetilde X\xrightarrow{\sim}\widetilde{f(X)}$. \unsure{Ainsi,|Aussi,}} l'on a bien
\[
 |X|^\circ=f(X)^\circ
\]
de sorte que (ii) implique également
\[
 X\,|\unsure{\circ|0}|\,Y\implies f(X)\,|\circ|\,f(Y).
\]
Il reste à vérifier que
\[
 X\mathrel{\overset{\circ}{\ll}}Y\implies f(X)\mathrel{\overset{\circ}{\ll}}f(Y),
\]
\deleted{Mais} va-t-on y arriver ? On aura une application \uline{croissante}
\[
 \unsure{\phi|\varphi}\mathbin{:}\widetilde X\longrightarrow\widetilde Y,\qquad \varphi(X)=Y
\]
\deleted{telle que} caractérisée par
\[
 X'\mathrel{\overset{\circ}{\ll}}\varphi(Y')
\]
et \deleted{\unsure{tout}} revient à ceci:
% slanted
\margin{cf \unsure{CF|CE}\,$\overline{\underline{\mathrm V}}$ p. \unsure{19}, 20, \unsure{Critère|critère} \unsure{b|3}}
\[
 X\mathrel{\overset{\circ}{\ll}}Y\implies |X|^\circ\subset |Y|^\circ
\]
\deleted{\unsure{On} sait déjà \unsure{q}} Par (i) on aura
\[
 |X|\subset |Y|\qquad\text{d'où}\quad |X|^\circ\subset |Y|,
\]
% archive sheet number: 69
